\documentclass[10pt,a4paper]{amsart}
\usepackage[margin=19mm]{geometry}
\usepackage{amsmath,amssymb,mathtools}
\usepackage[scr=rsfs,cal=euler]{mathalfa}
\usepackage{xfrac}
\usepackage[unicode,hidelinks,bookmarks=false]{hyperref}
\newcommand{\ie}{i.e.\ }
\newcommand{\cf}{cf.\ }
\newcommand{\resp}{respectively\ }
\newcommand{\Subsection}{\subsection}
\providecommand{\nobreakdash}{\nobreak\hbox{-}}
\setlength{\emergencystretch}{2em}
\multlinegap0pt
\usepackage{bm}
\usepackage{amssymb}
\usepackage{mathtools}
\newtheorem{theorem}{Theorem}
\title{Uncertain random geometric programming problems}
\author{Tapas Mondal, Akshay Kumar Ojha, Sabyasachi Pani}
\date{Source: arXiv:2310.01848v1 (CC BY 4.0)}
\begin{document}
\maketitle

\noindent\emph{Source and licence.} Uncertain random geometric programming problems. Version arXiv:2310.01848v1 (CC BY 4.0), distributed by arXiv under the Creative Commons Attribution 4.0 International licence, http://creativecommons.org/licenses/by/4.0/, as recorded in the arXiv licence metadata for this exact version and shown on the abstract page https://arxiv.org/abs/2310.01848v1. Full licence notice: \texttt{licenses/arxiv-2310.01848-CC-BY-4.0.txt}.\par\smallskip

\noindent\textit{Editorial scope.} Counted unit: the article's theorem thm.6 (source0.tex lines 463--473). The article's complete original proof of it is delivered with it (source0.tex lines 474--490, one \texttt{proof} environment of 17 lines). No mathematics is written, repaired or re-derived: the statement and the proof are the article's own lines, in the article's own notation, and no retained line is substituted or re-flowed.  No further source-local context is needed: every cross-reference of every retained line resolves inside the two delivered spans.  Nothing was omitted from any retained span, so the package carries no omission record.  No retained line contains a citation, so the wrapper carries no bibliography and no bibliography entry.  No retained line contains a figure environment, an image-inclusion command or drawing code, so no image is delivered and the package declares no asset dependency.  Editorial formatting only, in addition: an AMS \texttt{amsart} preamble that restates the article's own declarations and macros used by the retained lines, namely the environments \texttt{theorem} on separate counters, together with the standard packages those lines need.  The retained environments print in this document as Theorem 1 in order of appearance, whereas the article prints the same items with the numbering of its own full document.  The complete native source of the article as distributed in the arXiv e-print for this exact version is bundled unchanged as \texttt{sources/147927/source0.tex}.
\begin{theorem}\label{thm.6}
		If $\xi_l\sim\mathcal{N}(\mu_l,\sigma_l)$ with $l=1,2,\ldots,L$ are independent normal random variables, and $U_l$ with $l=1,2,\ldots,L$ are nonnegative variables, then for every tolerance level $\epsilon \in(0,0.5]$, 
		\begin{align*}
			Pr\bigg(\sum\limits_{l=1}^{L}\xi_lU_l\le 1\bigg)\ge 1-\epsilon
			\end{align*}
		is equivalent to
		\begin{align*}
			\sum\limits_{l=1}^{L}\mu_lU_l+\Phi^{-1}(1-\epsilon)\sqrt{\sum\limits_{l=1}^{L}\sigma_l^2U_l^2}\le 1,
		\end{align*}
	where $\Phi^{-1}(1-\epsilon)=\sqrt{2}\text{erf}^{-1}(1-2\epsilon)$ is the inverse of the standard normal distribution function.
	\end{theorem}

\begin{proof}
	Let $Z=\bigg(\sum\limits_{l=1}^{L}\xi_lU_l\bigg)-1.$ Then we have 
	$\mu_Z=E[Z]=E\bigg[\bigg(\sum\limits_{l=1}^{L}\xi_lU_l\bigg)-1\bigg]=\bigg(\sum\limits_{l=1}^{L}E[\xi_l]U_l\bigg)-1=\sum\limits_{l=1}^{L}\mu_lU_l-1$ and $\sigma_Z^2=var(Z)=var\bigg[\bigg(\sum\limits_{l=1}^{L}\xi_lU_l\bigg)-1\bigg]=\bigg(\sum\limits_{l=1}^{L}var[\xi_l]U_l^2\bigg)-0=\sum\limits_{l=1}^{L}\sigma_l^2U_l^2.$ Therefore, we have 
	\begin{equation*}
		\begin{aligned}
		Pr\bigg(\sum\limits_{l=1}^{L}\xi_lU_l\le 1\bigg)\ge 1-\epsilon
		 & \Leftrightarrow Pr(Z\le 0)\ge 1-\epsilon\\
	    & \Leftrightarrow Pr\bigg(\frac{Z-\mu_Z}{\sigma_Z}\le -\frac{\mu_Z}{\sigma_Z}\bigg)\ge 1-\epsilon\\
	    & \Leftrightarrow \Phi\bigg(-\frac{\mu_Z}{\sigma_Z}\bigg)\ge 1-\epsilon \qquad \bigg[\because \frac{Z-\mu_Z}{\sigma_Z}\sim \mathcal{N}(0,1)\bigg]\\ 
	    & \Leftrightarrow -\frac{\mu_Z}{\sigma_Z}\ge \Phi^{-1}(1-\epsilon)\\
	    & \Leftrightarrow \mu_Z+\sigma_Z \Phi^{-1}(1-\epsilon)\le0\\
	    & \Leftrightarrow \sum\limits_{l=1}^{L}\mu_lU_l-1+\Phi^{-1}(1-\epsilon)\sqrt{\sum\limits_{l=1}^{L}\sigma_l^2U_l^2}\le0\\
	    & \Leftrightarrow \sum\limits_{l=1}^{L}\mu_lU_l+\Phi^{-1}(1-\epsilon)\sqrt{\sum\limits_{l=1}^{L}\sigma_l^2U_l^2}\le1.
		\end{aligned}
	\end{equation*}
This completes the proof.
\end{proof}

\end{document}
